\documentclass[12pt,a4paper]{article}
\usepackage[utf8]{inputenc}
\usepackage[T1]{fontenc}
\usepackage{amsmath,amssymb}
\usepackage{amsthm}
\usepackage{hyperref}
\usepackage{geometry}
\geometry{margin=2.5cm}
\usepackage{enumitem}
\usepackage{booktabs}
\usepackage{xcolor}
\newtheorem{lemma}{Lemma}[section]
\newtheorem{theorem}{Theorem}[section]
\newtheorem{definition}{Definition}[section]

\title{\bfseries Formalising the Spectral Measure and Moment Identity \\ for the GDST Programme}
\author{Victor Geere \and Mentor}
\date{May 2026}

\begin{document}

\maketitle

\begin{abstract}
The Greedy Dirichlet Sub‑Sum Transform (GDST) proof of the Riemann Hypothesis relies on a 
moment identity linking a spectral measure $\mu_\sigma$ to the eigenvalues of an integral 
operator $T_\sigma$:
\[
\int_{\mathbb{R}} x^k \, d\mu_\sigma(x) = \sum_{i} \lambda_{\sigma,i}^k \qquad (k\ge 1,\; \sigma>0).
\]
At present this identity is assumed as an axiom. We outline a complete formalisation plan in 
Isabelle/HOL to construct $\mu_\sigma$ as a discrete measure (push‑forward of the counting 
measure on the index set under the eigenvalue map), to prove that integration with respect to 
it reduces to a sum over eigenvalues, and to connect this sum to the trace of $T_\sigma^k$ via 
its spectral decomposition. With this work the GDST theory will become significantly more 
self‑contained, removing the need for the external measure‑theoretic axiom.
\end{abstract}

\section{Motivation}
The GDST programme~\cite{Geere2026} develops a novel proof of the Riemann Hypothesis by 
associating a Dirichlet series $E(\theta,s)$ to a transfer operator $\mathcal{L}_s$, 
establishing a Fredholm determinant identity, and then building a spectral model (Jacobi matrix) 
whose spectrum is shown to consist of the ordinates of the non‑trivial zeros of $\zeta(s)$. 
A crucial step is the identity
\begin{equation}\label{eq:moment}
\int_{\mathbb{R}} x^k \, d\mu_\sigma(x) \;=\; \sum_{i} \lambda_{\sigma,i}^k
\;\;=\;\; \operatorname{tr}(T_\sigma^k) \qquad (\sigma>0,\;k\ge 1),
\end{equation}
where $\lambda_{\sigma,i}$ are the eigenvalues of the integral operator $T_\sigma$ (see 
Axioms~AX6a–c in the theory) and $\mu_\sigma$ is its spectral measure. In the current version 
of the formalisation, \eqref{eq:moment} is declared as an axiom.  To obtain a fully rigorous 
proof one must \emph{construct} $\mu_\sigma$ explicitly, prove the integration--summation 
equivalence, and then link the sum to the trace using the spectral theorem.

\section{Objectives}
\begin{enumerate}[label=(\arabic*)]
    \item Define a discrete spectral measure $\mu_\sigma$ as the push‑forward of the counting 
          measure on the index set (taken to be $\mathbb{N}$) under the eigenvalue map 
          $\lambda_\sigma : \mathbb{N} \to \mathbb{R}$.
    \item Prove that for any measurable $f:\mathbb{R}\to\mathbb{C}$ such that 
          $\sum_{i} |f(\lambda_{\sigma,i})| < \infty$,
          \[
          \int_{\mathbb{R}} f(x)\, d\mu_\sigma(x) \;=\; \sum_{i=0}^\infty f(\lambda_{\sigma,i}).
          \]
    \item Instantiate $f(x)=x^k$ to obtain the moment identity.
    \item Connect the eigenvalue sum to the trace of $T_\sigma^k$ via the spectral decomposition 
          (AX6a) and a minimal additional axiom (or ultimately a full proof).
    \item Replace the current axiom \texttt{AX\_moment} in the GDST theory by the proven theorem.
\end{enumerate}

\section{Prerequisites}
\begin{itemize}
    \item Isabelle/HOL with the \texttt{Analysis} library (measure theory, Lebesgue/Bochner 
          integral).
    \item The existing axiomatisation of the spectral decomposition of $T_\sigma$ (AX6a–c):
          \begin{quote}
          \texttt{T\_sigma\_spectral}: $T_\sigma f = \sum_i \lambda_{\sigma,i} \langle e_{\sigma,i}, f\rangle e_{\sigma,i}$,\\
          \texttt{lambda\_nonneg}: $\lambda_{\sigma,i} \ge 0$,\\
          \texttt{e\_orthonormal}: $(e_{\sigma,i})$ orthonormal in $L^2[0,\pi]$.
          \end{quote}
    \item The \texttt{Distr} and \texttt{count\_space} constructions from 
          \texttt{HOL-Analysis.Measure\_Notebook}.
    \item The \texttt{Trace\_Class} theory (imported in the GDST file).
\end{itemize}

\section{Phase 1: Construction of the Discrete Spectral Measure}
Let $I = \mathbb{N}$ be the index set.  The \textbf{counting measure} on $I$ is
$\kappa = \mathsf{count\_space}\; \mathbb{N}$, which assigns to a set $A\subseteq\mathbb{N}$ 
the number of elements (possibly $\infty$).  The eigenvalue function is 
$\lambda_{\sigma} : \mathbb{N} \to \mathbb{R}$, accessible as \texttt{lambda\_sigma $\sigma$}.
We define the spectral measure as the \textbf{distribution} of $\kappa$ under $\lambda_{\sigma}$:
\[
\mu_\sigma \;:=\; \mathrm{distr}\; \kappa\; \mathsf{borel}\; \lambda_{\sigma}.
\]
In Isabelle:
\begin{verbatim}
definition mu_sigma :: "real ⇒ real measure" where
  "mu_sigma σ = distr (count_space UNIV) borel (λi. lambda_sigma σ i)"
\end{verbatim}
Because $\kappa$ is $\sigma$-finite and $\lambda_\sigma$ is measurable (all functions from a 
discrete space are measurable), $\mu_\sigma$ is a well‑defined Borel measure on $\mathbb{R}$.  
Its closed support is the closure of $\{\lambda_{\sigma,i}\mid i\in\mathbb{N}\}$, which by 
AX6g will later be identified with the spectrum of the Jacobi matrix.

\paragraph{Multiplicity} If some eigenvalues coincide (i.e., $\lambda_{\sigma,i} = 
\lambda_{\sigma,j}$ for $i\neq j$), the push‑forward automatically adds the weights, so 
multiplicities are correctly represented.

\section{Phase 2: Integration with respect to the Spectral Measure}
Let $f:\mathbb{R}\to\mathbb{C}$ be measurable and assume the summability condition 
$\sum_{i=0}^\infty |f(\lambda_{\sigma,i})| < \infty$. Then $f$ is Bochner integrable with 
respect to $\mu_\sigma$ and
\[
\int_{\mathbb{R}} f \, d\mu_\sigma \;=\; \sum_{i=0}^\infty f(\lambda_{\sigma,i}).
\]
This follows from two standard lemmas:
\begin{lemma}[Integration of a push‑forward]
$\displaystyle \int_X f \, d(\mathrm{distr}\;M\; \mathcal{B}\; \phi) =
\int_Y (f \circ \phi) \, dM$, for $M$ a measure on $Y$ and $\phi:Y\to X$ measurable.
\end{lemma}
\begin{lemma}[Counting measure integral]
$\displaystyle \int_{\mathbb{N}} g \, d(\mathsf{count\_space}\; \mathbb{N}) =
\sum_{i=0}^\infty g(i)$ whenever the sum is absolutely summable (Bochner sense).
\end{lemma}
Applying the first lemma with $M = \kappa$, $\phi = \lambda_\sigma$, and then the second 
with $g = f \circ \lambda_\sigma$ yields the identity.  In Isabelle, the proof becomes:
\begin{verbatim}
lemma integral_mu_sigma:
  assumes "summable (λi. norm (f (lambda_sigma σ i)))"
  shows "∫ x. f x ∂(mu_sigma σ) = (∑ i. f (lambda_sigma σ i))"
  unfolding mu_sigma_def
  by (simp add: integral_distr integral_count_space_nat)
\end{verbatim}
(We may need to import the lemma \texttt{integral\_count\_space\_nat} from \texttt{HOL-Analysis}.)

\section{Phase 3: Moment Identity}
Take $f(x)=x^k$ with $k\ge 1$.  By the non‑negativity of eigenvalues (AX6b), the series 
$\sum_i \lambda_{\sigma,i}^k$ is absolutely summable provided the eigenvalues are, for instance, 
square‑summable (Hilbert–Schmidt condition).  In the GDST framework, $T_\sigma$ is trace‑class 
for $\sigma>0$ (this follows from the spectral decomposition and the super‑exponential growth 
of the kernel’s support, or can be taken as a mild additional assumption).  Thus 
$\sum_i \lambda_{\sigma,i} < \infty$, and consequently all higher integer moments converge.  

Applying Phase~2 we obtain:
\begin{lemma}[Moment identity]
For $\sigma>0$ and $k\ge 1$,
\[
\int_{\mathbb{R}} x^k \, d\mu_\sigma(x) \;=\; \sum_{i=0}^\infty (\lambda_{\sigma,i})^k .
\]
\end{lemma}
In Isabelle:
\begin{verbatim}
lemma moment_eq_sum:
  fixes σ k
  assumes "σ > 0" "k ≥ 1"
  shows "(∫ x. x ^ k ∂(mu_sigma σ)) = (∑ i. (lambda_sigma σ i) ^ k)"
  using integral_mu_sigma[of "λx. x^k"] summable_lambda_pow[of σ k] by auto
\end{verbatim}
where \texttt{summable\_lambda\_pow} is a lemma asserting the absolute summability of 
$\lambda_{\sigma,i}^k$, derivable from the trace‑class property.

\section{Phase 4: Connection to the Trace of $T_\sigma^k$}
The spectral theorem (AX6a) represents $T_\sigma$ as a multiplication operator.  For a 
compact self‑adjoint operator, the trace of $T_\sigma^k$ (when defined) equals the sum of the 
$k$-th powers of its eigenvalues:
\[
\operatorname{tr}(T_\sigma^k) \;=\; \sum_{i=0}^\infty (\lambda_{\sigma,i})^k .
\]
This is a classical fact but its full formalisation in the Hilbert space setting is involved.  
As a stepping stone, we temporarily adopt a minimal axiom:
\begin{verbatim}
axiomatization where
  trace_T_sigma_pow_eq_sum:
    "σ > 0 ⟹ trace (T_sigma σ ^^ k) = (∑ i. (lambda_sigma σ i) ^ k)"
\end{verbatim}
This axiom is \emph{independent} of the measure construction and can later be replaced by a 
complete proof using the trace definition in \texttt{Trace\_Class}.  Once available, the 
full moment‑trace identity follows:
\[
\int x^k d\mu_\sigma \;=\; \sum_i \lambda_{\sigma,i}^k \;=\; \operatorname{tr}(T_\sigma^k).
\]
This eliminates the former axiom \texttt{AX\_moment} and replaces it by a much milder (and 
more directly verifiable) statement.

\section{Phase 5: Integration into the GDST Theory}
\begin{itemize}[leftmargin=*]
    \item \textbf{Replace} the old \texttt{mu\_sigma} definition (a density on $[0,\pi]$) 
          with the new push‑forward definition.
    \item \textbf{Prove} the theorem \texttt{moment\_eq\_trace} as a combination of 
          \texttt{moment\_eq\_sum} and the trace axiom.
    \item \textbf{Remove} \texttt{AX\_moment} from the axiom registry.
    \item \textbf{Adjust} any subsequent lemmas that reference the old measure; the support 
          of the new measure is the closure of the eigenvalue set, which still satisfies 
          \texttt{J\_sigma\_spec} (AX6g). The $\sigma$‑independence argument now uses 
          the discrete nature of the measure, making the Hausdorff moment problem 
          application more direct.
\end{itemize}

\section{Timeline and Effort Estimate}
\begin{center}
\begin{tabular}{@{}ll@{}}
\toprule
\textbf{Phase} & \textbf{Estimated effort} \\
\midrule
1 \& 2: Discrete measure and integration theorem & 2--3 weeks \\
3: Moment identity (monomials) & 1 week \\
4: Trace axiom and linking & 1 week \\
5: Integration into GDST file & 1--2 weeks \\
\midrule
\textbf{Total} & \textbf{about 8 weeks (part‑time)} \\
\bottomrule
\end{tabular}
\end{center}
The work is routine measure theory and integration; the main Isabelle files to consult are 
\texttt{Distr.thy}, \texttt{Bochner\_Integration.thy}, and \texttt{Countable.thy}.  
Familiarity with the Lebesgue integral on discrete spaces is required.

\section{Conclusion}
This programme transforms an opaque axiom into a transparent, formally verified identity 
resting only on the spectral decomposition of $T_\sigma$ and the standard relationship between 
eigenvalue sums and operator traces (the latter itself reducible to the spectral theorem).  
The formalisation is modular: the measure‑theoretic part is completely self‑contained and 
can be developed independently of the deeper GDST structure.  Upon completion the GDST proof 
of the Riemann Hypothesis will have one fewer external assumption, moving closer to a fully 
self‑contained formal verification.

\begin{thebibliography}{9}
\bibitem{Geere2026}
V.~Geere, \emph{GDST Programme: Proof of the Riemann Hypothesis via the Greedy Dirichlet 
Sub‑Sum Transform}, Isabelle theory file, 2026.
\bibitem{IsabelleAnalysis}
T.~H\"olzl et al., \emph{Isabelle/HOL Analysis Library}, AFP.
\bibitem{Simon2011}
B.~Simon, \emph{Szegő's Theorem and its Descendants}, Princeton Univ. Press, 2011.
\end{thebibliography}

\end{document}